% !TeX root = ../papers/nc-p3.tex
\section{分岐を通る短い周期点加群と無限生成イデアル}
\label{sec:periodic-branch-obstruction}
\begin{theorem}\label{thm:periodic-branch-obstruction}
曲面射影のない標準正則窓に片側直線ファイバーがあれば、四周期を法として
順側分岐中心は高々20個であり、そのいずれかを通る長さ40以下の周期点列がある。
逆側でも同じである。その右周期点加群と逆向きの左周期点加群の表示イデアルは
それぞれ無限生成である。従ってこの残余枝が実在すれば右・左の両方でnoetherian性に失敗する。
\end{theorem}
\begin{proof}
各辺で六双一次式の零点集合の次元は高々1である。実際、曲線上で動く直線ファイバーはなく、
各孤立中心上のファイバーも高々直線である。この線型系から一般の五式を取れば
$\mathbb P^3\times\mathbb P^3$ で曲線の完全交叉となる。各段で正次元の余分な成分を次の式で切れ、
基底集合が次元1以下なので余次元5まで正則列になるからである。
完全交叉の類 $(H_1+H_2)^5$ により二つの射影方向の次数は10ずつである。
全ての点対応の既約曲線成分はこの交叉の成分なので、両次数の和の総計は一辺で20以下である。
定理\ref{thm:finite-branch-corridors}の各列は、辺ごとに異なる曲線成分を用いる。
一般前点を逆にたどって最初の順側中心を復元できるので、異なる列の間でも重複しない。
列の辺の両次数の和は $2\sum_{j=1}^m d_j$ であるから、
四位置の全列について $\sum m\le\sum d_j\le40$ である。各列は $m\ge2$ なので中心数の上界を得る。
各中心からその列の内部で到達できる順側中心を一つ選ぶ。この有限有向グラフの単純巡回路に沿って
元の点経路を連結すると、その長さは巡回路上の各列の $m$ の和以下である。
この長さは四位置を元に戻す。左向きの結論は語順の反転による。

周期点列 $(p_j)$ の一つの始点 $i$ からの右点加群を $M$ とし、
$M_{i+n}=k$ に次の文字を $p_{i+n}$ で評価して作用させる。
終点 $j=i+n$ の左単純のKoszul解消を $M$ とテンソルする。
$n\ge2$ の第一ホモロジーは
\[
 \operatorname{Tor}^A_1(M,k_j)
   =\ker(p_{j-1}:V_{j-1}\to k)\,/\,R_{j-2}(p_{j-2}) .
\]
実際これは $R_{j-2}\to V_{j-1}\to k$ の二微分そのものである。
従ってその次元は順側中心を訪れた二文字後で1、それ以外は0である。
$n=1$ では一次核が三次元である。極小生成元と第一Torの同定はイデアルの有限生成性を要しない。
各周期に一回以上新たな生成元が必要であるから、$e_i A\to M$ の核は無限生成である。
逆周期列を使う左側も同様である。
\end{proof}

\begin{proposition}\label{prop:terminal-branch-contact}
順側中心を頂点とし、各中心 $p$ の曲線列で $x_j(t)$ が順側中心になる全ての場合に、
$p$ から $x_j(t)$ への辺を置く。この有限有向グラフの終端強連結成分内の各 $p$ について、
その列の順側階数落ちはただ一つのパラメータ $t$ に集中し、到達する各 $x_j$ はそこで不分岐である。
特に長さ2の曲線列がこの成分から出発するとき、次の順側中心の階数二スキームは非被約である。
\end{proposition}
\begin{proof}
異なる二つのパラメータから成分内の中心に着けば、それぞれに $p$ への帰路を付けられる。
二つの閉経路は最初の直線で異なる点を選んでいる。反復で長さを揃えた後、二つを任意順に選べば
指数個の独立な経路評価が生じる（定理\ref{thm:recurrent-point-correspondences}の二分岐の証明）。
また $x_j$ の微分が零なら、そのパラメータの双対数近傍は $p$ と $x_j(t)$ を固定しながら
最初の直線の点を非自明に動かす。定数の帰路を付け、独立な双対数を毎回使えば
定理\ref{thm:no-two-sided-centers}の最後と同じく指数個の独立評価を得る。
$m=2$ なら $d_2=1$、従って $f_1=3-d_2=2$ である。一つの点で二次の接触をするため、
その順側階数二スキームは $x_1$ の非零接ベクトルを含む。
中心は孤立しているのでこのスキームはそこで非被約である。
\end{proof}
